
\documentclass[a4paper]{article} %

\date{} %

% Please, do not change any of the above lines

\usepackage{amssymb} % add other LaTeX packages you need here.

\begin{document}
\setcounter{page}{1} % Please, do not delete this line

% Your title
\title{Solving lattice-valued relational equations and inequations}

% Authors
\author{Branimir \v{S}e\v{s}elja,  \\ %
       University of Novi Sad\\ \\ % Affiliation 1
       Vanja Stepanovi\'{c} \\ % If any other author with different Affilation
       University of Belgrade\\ \\ % Affiliation 2 (if needed)
       Andreja Tepav\v{c}evi\'{c} \\
       Mathematical Institute SANU, Belgrade, Serbia \\
       % Add authors and affiliation as needed
       \texttt{andreja@dmi.uns.ac.rs} % Only one corresponding e-mail, please
       }%

\maketitle
%%%%%%

We will present our investigation on the existence of optimal solutions to lattice‑valued relational equations and inequations, extending the classical notions of greatest and least solutions to maximal and minimal ones \cite{acta, comsis, ST, 2026}. A key result establishes that infinite distributivity of infimum relative to supremum in the codomain lattice is both necessary and sufficient for typical inequations to admit a greatest solution. We show that minimal solutions may fail to exist even under strong lattice assumptions, while meet‑continuity, though sufficient for maximal solutions, is not always necessary.

A distinctive aspect of our work is the connection between approximate solutions and factor algebras, where weak congruences provide the structural framework for transferring solvability properties \cite{gjba}. This perspective links the study of relational systems with algebraic structures, clarifying how approximate solutions can be systematically understood through weak-congruence theory.

Connected to this are quotient structures of lattice-valued matrix spaces, constructive algorithms for least solutions that terminate after countably many steps, and refined lattice‑theoretic conditions for solution existence. Altogether, the results advance a connection of the theory of lattice‑valued equations and inequations with the theory of weak congruences (solutions of equations on factor subalgebras).

\vspace{0.5cm}
This 
 research was supported by the Science Fund of the Republic of Serbia, \# Grant no 6565,  Advanced Techniques of Mathematical Aggregation and Approximative Equations Solving in Digital Operational Research-AT-MATADOR. 


%%%%%%



\begin{thebibliography}{99}

 \bibitem{gjba} M.Z.\ Grulovi\' c, J.\ Jovanovi\' c, B.\ \v{S}e\v{s}elja, A.\ Tepav\v{c}evi\'c,   \emph{Lattice Characterization of Some Classes of Groups by Series of Subgroups},  Int. J. Algebra Comput. 33 (02), 211-235 (2023).
 
\bibitem{acta} Stepanovi\'c, V., Tepav\v{c}evi\'c: A. A Note on the Solutions to Lattice-valued Relational Equations and Inequations, Acta Polytechnica Hungarica \textbf{22}(12), 27--46, 2025.

\bibitem{comsis} Stepanovi\'c, V., Tepav\v{c}evi\'c: Complete meet-continuous codomain lattice and extremal solutions of L-valued fuzzy relation equations, accepted for publication in \emph{Computer Science and Information Systems}, 2026.


\bibitem{ST} Stepanovi\'c, V., Tepav\v{c}evi\'c, A.: Fuzzy sets (in)equations with a complete codomain lattice. Kybernetika \textbf{58}, 145--162, 2022.% \url{https://doi.org/10.14736/kyb-2022-2-0145}.

\bibitem{2026} Stepanovi\'c, V., Tepav\v{c}evi\'c, A.:
Systems of fuzzy relational equations in the general lattice-valued case,
International Journal of Approximate Reasoning
\textbf{195},
109703, 2026. 


\end{thebibliography}

\end{document}

